\chapter{Pochodna Frécheta i różniczka zupełna}
\label{ma2:ch:frechet}

Od tego rozdziału stosujemy ogólne twierdzenia rachunku
różniczkowego w rzeczywistych przestrzeniach unormowanych.
Zupełność zakładamy tylko tam, gdzie jest potrzebna.
Podstawowe wyniki pochodzą z analizy funkcjonalnej,
rozdział 14 (AF 14); wykorzystujemy ich dowody, a w razie
osłabienia hipotez wskazujemy uzasadnienie.
Przypadek euklidesowy pozostaje szczególnie ważny z uwagi
na rachunek współrzędnych i geometrię.

\section{Liniowe przybliżenie przyrostu}

Niech $E,F$ będą rzeczywistymi przestrzeniami unormowanymi,
$U\subseteq E$ otwarty i $p\in U$.
Przestrzeń operatorów liniowych ciągłych
$\mathcal L(E,F)$ wyposażamy w normę operatorową
(stwierdzenie \ref{prop:operator-norm12}).

\begin{definicja}[Pochodna Frécheta]
\label{ma2:def:frechet}
Odwzorowanie $f:U\to F$ nazywamy różniczkowalnym
w sensie Frécheta w $p$, jeśli istnieje
$A\in\mathcal L(E,F)$ takie, że
\[
 f(p+h)=f(p)+Ah+r(h),\qquad
 \lim_{\substack{h\to0\\h\ne0}}
 \frac{\|r(h)\|_F}{\|h\|_E}=0.
\]
Reszta jest określona dla $h$ z pewnej kuli wokół zera
i $r(0)=0$. Piszemy $r(h)=o(\|h\|_E)$.
Operator $A$ oznaczamy przez $Df(p)$, a także
$\mathrm df_p$ (różniczka zupełna w $p$).
Jest to definicja AF 14.3.
\end{definicja}

Dla $E=\{0\}$ warunek na resztę rozumiemy w sensie
$\varepsilon$--$\delta$ (warunek dla niezerowych $h$ jest pusty); jedyny operator
liniowy jest wtedy zerowy.
Dla $E=F=\R$ otrzymujemy definicję pochodnej i różniczki
z AM I 20.1.1--2 oraz 20.2.1, gdyż każdy operator liniowy ma
postać $h\mapsto ah$ i iloraz reszty przez $h$ dąży
do zera dokładnie wtedy, gdy jej moduł podzielony
przez $|h|$ dąży do zera.

W ogólnych przestrzeniach unormowanych pochodną
kierunkową definiujemy tym samym ilorazem co w definicji
\ref{ma2:def:directional}, z granicą w normie $F$.
Pochodna Gâteaux wymaga dodatkowo ograniczoności
operatora liniowego; w nieskończonym wymiarze sama
liniowość nie wystarcza do jego ciągłości (AF 14.2).

\begin{twierdzenie}[Ciągłość i pochodna Gâteaux]
\label{ma2:thm:frechet-gateaux}
Jeśli $f$ jest różniczkowalne w sensie Frécheta w $p$,
to jest ciągłe w $p$. Dla każdego $v\in E$ istnieje
dwustronna pochodna kierunkowa i
$\partial_v f(p)=Df(p)v$.
W szczególności pochodna Frécheta jest jednoznaczna
i $D_Gf(p)=Df(p)$, gdzie pochodna Gâteaux oznacza
ciągły operator liniowy reprezentujący wszystkie
pochodne kierunkowe (AF 14.2).
\end{twierdzenie}

Jest to dokładnie twierdzenie AF 14.5
(\texttt{thm:frechet-implies-gateaux14}); tam znajduje
się dowód dla przestrzeni unormowanych. W szczególności
nie wymaga on zupełności. Dla $E=\R^n$, $F=\R^m$
zastosowanie stwierdzenia \ref{ma2:prop:gateaux-matrix}
daje macierz $J_f(p)$ operatora $Df(p)$.

\begin{stwierdzenie}[Jednostajność względem kierunku]
\label{ma2:prop:frechet-uniform}
Jeśli $E\ne\{0\}$ i $A\in\mathcal L(E,F)$, to
$Df(p)=A$ wtedy i tylko wtedy, gdy
\[
 \lim_{t\to0,\ t\ne0}
 \sup_{\|v\|_E=1}
 \left\|\frac{f(p+tv)-f(p)}{t}-Av\right\|_F=0.
\]
Parametr $t$ jest rzeczywisty i ma moduł dostatecznie
mały, aby $p+tv\in U$ dla wszystkich jednostkowych $v$.
Supremum może a priori przyjąć wartość $+\infty$.
\end{stwierdzenie}

Stosujemy AF 14.6 (\texttt{prop:uniform-directions14}).
W porównaniu z definicją pochodnej Gâteaux nie wystarcza
granica dla każdego ustalonego $v$: potrzebna jest
jedna kontrola błędu dla wszystkich jednostkowych
kierunków (por. twierdzenie \ref{ma2:thm:radial-limit}).

\begin{przyklad}[Ciągłość i pochodna Gâteaux nie wystarczają]
\label{ma2:ex:frechet-counterexample}
W przykładzie AF 14.9 funkcja $f:\R^2\to\R$ określona
przez $f(0,0)=0$ oraz
$f(x,y)=x^3y/(x^4+y^2)$ poza zerem jest ciągła w zerze
i $D_Gf(0,0)=0$. Na paraboli $y=x^2$ dla $x\ne0$ mamy jednak
\[
 \frac{|f(x,x^2)|}{\|(x,x^2)\|_2}
 =\frac{1}{2\sqrt{1+x^2}}\longrightarrow\frac12.
\]
Pochodna Frécheta, gdyby istniała, musiałaby być zerowa
na mocy twierdzenia \ref{ma2:thm:frechet-gateaux},
co przeczy warunkowi na resztę. Pełne obliczenia
ciągłości i pochodnych kierunkowych są w AF 14.9.
\end{przyklad}

\section{Reguły różniczkowania}

\begin{twierdzenie}[Kombinacje liniowe i reguła łańcuchowa]
\label{ma2:thm:frechet-chain}
Niech $E,F,G$ będą rzeczywistymi przestrzeniami
unormowanymi, $U\subseteq E$, $V\subseteq F$ otwarte,
$p\in U$. Jeżeli $f_1,f_2:U\to F$ są różniczkowalne
w $p$, to dla $\alpha,\beta\in\R$ zachodzi
\[
 D(\alpha f_1+\beta f_2)(p)=
 \alpha Df_1(p)+\beta Df_2(p).
\]
Jeżeli $f:U\to V$ jest różniczkowalne w $p$,
a $g:V\to G$ w $f(p)$, to
\[
 D(g\circ f)(p)=Dg(f(p))\circ Df(p).
\]
Odwzorowanie afiniczne $x\mapsto Lx+b$ dla
$L\in\mathcal L(E,F)$, $b\in F$ ma pochodną $L$.
\end{twierdzenie}
\begin{proof}
Korzystamy ze stwierdzenia AF 14.14 oraz twierdzenia
AF 14.15. W źródle założono zupełność przestrzeni,
lecz dowody jej nie wykorzystują: w pierwszym
sumuje się reszty i używa nierówności trójkąta,
a w drugim oznacza się $A=Df(p)$, $B=Dg(f(p))$,
$k=f(p+h)-f(p)=Ah+r_f(h)$. Dla małych $h$ mamy
$\|r_f(h)\|\le\|h\|$, stąd
$\|k\|\le(\|A\|+1)\|h\|$.
Reszta po odjęciu $BAh$ od przyrostu złożenia
to $Br_f(h)+r_g(k)$. Pierwszy składnik ma normę
co najwyżej $\|B\|\|r_f(h)\|=o(\|h\|)$.
Dla $k\ne0$ zachodzi
\[
 \frac{\|r_g(k)\|}{\|h\|}
 \le(\|A\|+1)\frac{\|r_g(k)\|}{\|k\|}\longrightarrow0,
\]
a dla $k=0$ reszta $r_g(k)$ jest zerowa.
Operator $BA$ jest liniowy i ograniczony przez
$\|B\|\|A\|$.
Wszystkie działania odbywają się na istniejących
wektorach i ograniczonych operatorach, bez przejścia
do granicy ciągu Cauchy'ego. W przypadku afinicznym
reszta jest identycznie zerowa. Te same dowody
obowiązują zatem dla przestrzeni unormowanych.
\end{proof}

\begin{uwaga}[Dlaczego pochodna Gâteaux nie wystarcza?]
W AF 14.11 (stwierdzenie
\texttt{prop:chain-fails14}) do funkcji $f$ z przykładu
\ref{ma2:ex:frechet-counterexample} stosujemy krzywą
$\gamma(t)=(t,t^2)$. Wtedy $(f\circ\gamma)(t)=t/2$
dla każdego $t$, więc $(f\circ\gamma)'(0)=1/2$,
podczas gdy $D_Gf(0,0)\gamma'(0)=0$.
Krzywa $\gamma$ jest nawet różniczkowalna w sensie
Frécheta, gdyż $\gamma(t+s)-\gamma(t)=(s,2ts+s^2)$.
Oznacza to, że reguła łańcuchowa może zawieść już przy
zastąpieniu pochodnej zewnętrznego odwzorowania
samą pochodną Gâteaux.
\end{uwaga}

\begin{twierdzenie}[Reguła dla odwzorowań dwuliniowych]
\label{ma2:thm:frechet-bilinear}
Niech $E,E_1,E_2,G$ będą rzeczywistymi przestrzeniami
unormowanymi i $B:E_1\times E_2\to G$ dwuliniowym
odwzorowaniem spełniającym
$\|B(u,v)\|\le C\|u\|\|v\|$ dla pewnego $C\ge0$.
Wówczas $B$ jest różniczkowalne na całym iloczynie
$E_1\times E_2$ wyposażonym w normę maksimum i
\[
 DB(u,v)(h,k)=B(h,v)+B(u,k).
\]
Jeżeli $f:U\to E_1$ i $g:U\to E_2$ są różniczkowalne
w $p$, gdzie $U\subseteq E$ jest otwarty i $p\in U$, to
$z\mapsto B(f(z),g(z))$ jest różniczkowalne w $p$ i
\[
 D(B(f,g))(p)h
 =B(Df(p)h,g(p))+B(f(p),Dg(p)h).
\]
\end{twierdzenie}
\begin{proof}
Dowód AF 14.16 jest ważny bez założenia zupełności:
w normie maksimum produktu reszta jest równa
$B(h,k)$ i ma normę co najwyżej
$C\max(\|h\|,\|k\|)^2$.
Wskazany operator liniowy jest ograniczony, ponieważ
$\|B(h,v)+B(u,k)\|\le C(\|u\|+\|v\|)
\max(\|h\|,\|k\|)$.
W dowodzie drugiej części pochodna pary $(f,g)$
jest operatorem $h\mapsto(Df(p)h,Dg(p)h)$
o normie nie większej niż $\max(\|Df(p)\|,\|Dg(p)\|)$.
Iloraz normy reszty przez $\|h\|$ jest maksimum
dwóch ilorazów dążących do zera. Stosujemy następnie
twierdzenie \ref{ma2:thm:frechet-chain}.
\end{proof}

Dla $E_1=E_2=G=\R$ oraz $B(u,v)=uv$ otrzymujemy
regułę iloczynu. Niech $f,g:U\to\R$ będą
różniczkowalne w $p\in U$ i $g(p)\ne0$. Wówczas ciągłość $g$ w $p$ zapewnia,
że mianownik nie znika w pewnej kuli wokół $p$.
Stosując regułę łańcuchową do $s\mapsto1/s$
(AM I 20.3.2), a następnie regułę iloczynu, dostajemy
\[
 D(f/g)(p)h
 =\frac{g(p)Df(p)h-f(p)Dg(p)h}{g(p)^2}.
\]

\begin{przyklad}[Nieskończony wymiar]
Niech $E=C([0,1],\R)$ z normą supremum
(por. twierdzenie \ref{ma2:thm:continuous}).
Odwzorowanie $P:E\to E$, $P(u)=u^2$, jest
różniczkowalne, a $DP(u)h=2uh$.
Istotnie, mnożenie punktowe jest dwuliniowe
i $\|uv\|_\infty\le\|u\|_\infty\|v\|_\infty$.
Twierdzenie \ref{ma2:thm:frechet-bilinear} zastosowane
do $B(u,u)$ daje wskazany operator. Reszta $h^2$
ma normę $\|h\|_\infty^2$.
\end{przyklad}

\section{Pochodne częściowe na iloczynach}

Niech $E=E_1\times\cdots\times E_N$, gdzie $N\ge1$,
a $E_j$ są rzeczywistymi przestrzeniami unormowanymi.
Stosujemy normę $\|h\|=\max_j\|h_j\|$.
Przez $i_j:E_j\to E$ oznaczamy włożenie wektora
na $j$-tą pozycję, z zerami na pozostałych.
Niech $F$ będzie rzeczywistą przestrzenią unormowaną,
$U\subseteq E$ otwarty i $f:U\to F$.
Dla każdego $j=1,\ldots,N$ oraz $p\in U$ definiujemy $D_jf(p)$ jako pochodną
Frécheta funkcji $z\mapsto f(p+i_jz)$ w $z=0$,
jeśli istnieje. Dziedzina tej funkcji,
$\{z\in E_j:p+i_jz\in U\}$, jest otwarta i zawiera zero,
gdyż $i_j$ jest liniową izometrią.
Jest to pochodna częściowa względem całego bloku
zmiennych, a nie pojedynczej współrzędnej skalarnej.

\begin{stwierdzenie}
\label{ma2:prop:frechet-blocks}
Jeżeli $Df(p)$ istnieje, to istnieją wszystkie
$D_jf(p)\in\mathcal L(E_j,F)$ i
\[
 D_jf(p)=Df(p)\circ i_j,\qquad
 Df(p)h=\sum_{j=1}^N D_jf(p)h_j.
\]
\end{stwierdzenie}
\begin{proof}
Odwzorowanie $z\mapsto p+i_jz$ jest afiniczne
i jego pochodną jest $i_j$ (operator ograniczony,
gdyż $\|i_jz\|=\|z\|$). Reguła łańcuchowa daje
pierwszą równość. Druga wynika z
$h=\sum_j i_jh_j$ i liniowości $Df(p)$.
\end{proof}

\begin{twierdzenie}[Kryterium ciągłych pochodnych częściowych]
\label{ma2:thm:continuous-block-partials}
Niech $F$ będzie rzeczywistą przestrzenią unormowaną,
$U\subseteq E$ zbiorem otwartym, $f:U\to F$ i $p\in U$.
Jeżeli wszystkie pochodne częściowe $D_jf$ istnieją
w pewnym otoczeniu $p$ i są ciągłe w $p$ w normach
operatorowych, to $f$ jest różniczkowalne w sensie
Frécheta w $p$ i
\[
 Df(p)h=\sum_{j=1}^N D_jf(p)h_j.
\]
\end{twierdzenie}
\begin{proof}
Oznaczmy $A_j=D_jf(p)$ i
$Ah=\sum_j A_jh_j$. Operator $A$ jest liniowy
i ograniczony, bo $\|Ah\|\le(\sum_j\|A_j\|)\|h\|$.
Dla $\varepsilon>0$ wybieramy $\delta>0$ takie,
że kula $B(p,\delta)$ zawiera się w otoczeniu
z założenia i dla $z$ w tej kuli oraz każdego $j$
zachodzi $\|D_jf(z)-A_j\|<\varepsilon/N$.
Można użyć wspólnego promienia, bo bloków jest skończenie wiele.

Niech $\|h\|<\delta$, $p_0=p$ oraz
$p_j=p+\sum_{k=1}^j i_k h_k$.
Odcinek $p_{j-1}+t i_jh_j$, $0\le t\le1$,
zawiera się w $B(p,\delta)$, ponieważ w normie
maksimum odległość od $p$ nie przekracza $\|h\|$.
Funkcja
\[
 \gamma_j(t)=f(p_{j-1}+t i_jh_j)-t A_jh_j
\]
jest różniczkowalna (włącznie z pochodnymi jednostronnymi
na końcach): z definicji pochodnej częściowej,
po przyroście parametru $s$ przyrost $j$-tego bloku
wynosi $s h_j$. Jeśli $h_j=0$, krzywa jest stała.
W przeciwnym przypadku, dla $z=p_{j-1}+t i_jh_j$,
z definicji pochodnej częściowej mamy
$f(z+i_jk)-f(z)=D_jf(z)k+r_{j,z}(k)$,
gdzie $\|r_{j,z}(k)\|=o(\|k\|)$ dla $k\to0$.
Dla $k=sh_j$ otrzymujemy
\[
 \frac{\|r_{j,z}(sh_j)\|}{|s|}
 =\|h_j\|\frac{\|r_{j,z}(sh_j)\|}{\|sh_j\|}
 \longrightarrow0.
\]
Daje to w obu przypadkach
\[
 \gamma_j'(t)=
 (D_jf(p_{j-1}+t i_jh_j)-A_j)h_j,
 \qquad
 \|\gamma_j'(t)\|\le\frac{\varepsilon}{N}\|h_j\|.
\]
Różniczkowalność zapewnia ciągłość $\gamma_j$ na odcinku.
Lemat \ref{ma2:lem:vector-increment}, który nie wymaga
zupełności przestrzeni wartości, daje
\[
 \|f(p_j)-f(p_{j-1})-A_jh_j\|
 \le\frac{\varepsilon}{N}\|h_j\|.
\]
Sumując po $j$ i używając sumy teleskopowej, mamy
\[
 \|f(p+h)-f(p)-Ah\|
 \le\frac{\varepsilon}{N}\sum_{j=1}^N\|h_j\|
 \le\varepsilon\|h\|.
\]
Ponieważ $\varepsilon$ było dowolne, jest to
warunek różniczkowalności z pochodną $A$.
\end{proof}

\section{Klasa \texorpdfstring{$C^1$}{C1} i rachunek współrzędnych}

\begin{definicja}
\label{ma2:def:c1}
Piszemy $f\in C^1(U,F)$, jeśli $f$ jest różniczkowalne
w sensie Frécheta w każdym punkcie $U$ oraz
$Df:U\to\mathcal L(E,F)$ jest ciągłe w normie operatorowej.
Jest to definicja AF 14.18 rozszerzona na przestrzenie
unormowane; zupełność nie jest konieczna do jej
sformułowania.
\end{definicja}

\begin{wniosek}
\label{ma2:cor:c1-blocks}
Niech $U$ będzie otwartym podzbiorem skończonego
iloczynu $E$ rzeczywistych przestrzeni unormowanych
i $f:U\to F$, gdzie $F$ jest rzeczywistą przestrzenią
unormowaną. Wówczas $f$ jest klasy $C^1$ wtedy i tylko wtedy, gdy wszystkie
$D_jf$ istnieją na $U$ i są ciągłe w normach operatorowych.
\end{wniosek}
\begin{proof}
Jeśli $f$ jest klasy $C^1$, to
$D_jf(p)=Df(p)\circ i_j$ i
$\|D_jf(p)-D_jf(q)\|\le\|Df(p)-Df(q)\|$.
Odwrotnie, twierdzenie
\ref{ma2:thm:continuous-block-partials} daje pochodną
w każdym punkcie, a z jej wzoru wynika
$\|Df(p)-Df(q)\|\le\sum_j\|D_jf(p)-D_jf(q)\|$.
Skończona suma dąży do zera przy $q\to p$.
\end{proof}

\begin{stwierdzenie}[Niezależność od wyboru norm równoważnych]
\label{ma2:prop:frechet-equivalent-norms}
Zastąpienie norm w $E,F$ normami równoważnymi nie
zmienia istnienia ani wartości pochodnej Frécheta,
ani przynależności odwzorowania do klasy $C^1$.
\end{stwierdzenie}
\begin{proof}
Niech $a_E\|h\|\le\|h\|_*\le b_E\|h\|$
i $a_F\|y\|\le\|y\|_*\le b_F\|y\|$,
gdzie wszystkie cztery stałe są dodatnie.
Otwartość dziedziny nie zmienia się na mocy twierdzenia
\ref{ma2:thm:equiv}. Dla dowolnego liniowego $T:E\to F$ i
$r(h)=f(p+h)-f(p)-Th$ oraz $h\ne0$ mamy
\[
 \frac{a_F}{b_E}\frac{\|r(h)\|}{\|h\|}
 \le\frac{\|r(h)\|_*}{\|h\|_*}
 \le\frac{b_F}{a_E}\frac{\|r(h)\|}{\|h\|}.
\]
Stosując te same nierówności do $Th$ zamiast $r(h)$,
widzimy najpierw, że operator liniowy jest ograniczony
w normach pierwotnych wtedy i tylko wtedy, gdy jest
ograniczony w nowych normach. Identyczne stałe dają
$(a_F/b_E)\|T\|\le\|T\|_*\le(b_F/a_E)\|T\|$,
korzystając z równoważnego wzoru na normę operatorową
jako supremum $\|Th\|/\|h\|$ po $h\ne0$.
Dla przestrzeni zerowej operator jest zerowy i
nierówności pozostają prawdziwe.
Zatem ograniczoność operatora, zbieżność ilorazu
reszty do zera i ciągłość pochodnej są niezależne
od wyboru tych norm.
\end{proof}

W szczególności dla $E=\R^n$, $F=\R^m$
istnienie i ciągłość wszystkich pochodnych cząstkowych
funkcji współrzędnych na $U$ są równoważne
przynależności do klasy $C^1$.
Istotnie, dla $E_j=\R$ istnienie pochodnej cząstkowej
$\partial_j f(p)$ jest równoważne istnieniu pochodnej
Frécheta przekroju jednej zmiennej, gdyż
\[
 \left\|\frac{f(p+te_j)-f(p)}{t}-\partial_j f(p)\right\|
 =\frac{\|f(p+te_j)-f(p)-t\partial_j f(p)\|}{|t|}.
\]
Wtedy $D_jf(p)t=t\partial_j f(p)$, a norma różnicy
dwóch takich operatorów jest równa normie różnicy
ich wektorów współczynników. Istnienie i ciągłość
wektora pochodnej cząstkowej sprowadzają się do
odpowiednich własności jego współrzędnych
(twierdzenie \ref{ma2:thm:multi-limit} zastosowane
do ilorazu różnicowego oraz do granicy wektora).
Operator z twierdzenia
\ref{ma2:thm:continuous-block-partials} ma macierz
Jacobiego. Zmiana normy maksimum na euklidesową
jest dozwolona na mocy stwierdzenia
\ref{ma2:prop:frechet-equivalent-norms} i twierdzenia
\ref{ma2:thm:finite}.

\begin{wniosek}[Złożenie odwzorowań klasy $C^1$]
\label{ma2:cor:c1-chain}
Niech $E,F,G$ będą rzeczywistymi przestrzeniami
unormowanymi, $U\subseteq E$, $V\subseteq F$ otwarte.
Jeżeli $f\in C^1(U,F)$, $f(U)\subseteq V$ oraz
$g\in C^1(V,G)$, to $g\circ f\in C^1(U,G)$.
\end{wniosek}
\begin{proof}
Z twierdzenia \ref{ma2:thm:frechet-chain} pochodna
złożenia istnieje w każdym punkcie i wynosi
$Dg(f(x))Df(x)$. Dla ustalonego $p\in U$ zachodzi
\begin{align*}
 &\|Dg(f(x))Df(x)-Dg(f(p))Df(p)\|\\
 &\quad\le \|Dg(f(x))\|\|Df(x)-Df(p)\|\\
 &\qquad+\|Dg(f(x))-Dg(f(p))\|\|Df(p)\|.
\end{align*}
Ponieważ $f$ jest ciągłe, $Dg(f(x))\to Dg(f(p))$,
więc normy $\|Dg(f(x))\|$ są ograniczone w pewnym
otoczeniu $p$. Ponadto $Df(x)\to Df(p)$. Oba składniki
prawej strony dążą do zera, co daje ciągłość pochodnej.
\end{proof}

Istnienie pochodnej w każdym punkcie nie daje jeszcze
klasy $C^1$. Przykład AF 14.22 pokazuje to już na
$\R$: $f(0)=0$, $f(x)=x^2\sin(1/x^2)$ dla $x\ne0$.
Pochodna w zerze istnieje i wynosi zero, ale pochodna
nie jest ciągła w zerze (pełne obliczenia w AF 14.22).

Dla $f:\R^n\supseteq U\to\R^m$ różniczkowalnego w $p$
różniczka ma postać
\[
 \mathrm df_p=\sum_{j=1}^n\partial_j f(p)\,\mathrm dx_j,
 \qquad \mathrm dx_j(h)=h_j.
\]
Iloczyn wektora i funkcjonału oznacza tu operator
$h\mapsto h_j\partial_j f(p)$. Dla funkcji skalarnej
otrzymujemy $\mathrm df_p(h)=\langle\nabla f(p),h\rangle$.
W rzeczywistej przestrzeni Hilberta reprezentację
$Df(p)h=\langle h,\nabla f(p)\rangle$ zapewnia
twierdzenie Riesza (AF 14.24); w ogólnej przestrzeni
unormowanej różniczka funkcji skalarnej jest funkcjonałem
liniowym ciągłym, nie zaś kanonicznie wektorem.

\begin{przyklad}
Niech $f(x,y)=(x^2+y,xy)$ i $g(u,v)=u+v^2$.
Wszystkie pochodne cząstkowe są ciągłe, więc
$f,g$ są klasy $C^1$ na odpowiednich przestrzeniach.
Z reguły łańcuchowej
\[
 J_{g\circ f}(x,y)
 =\begin{pmatrix}1&2xy\end{pmatrix}
  \begin{pmatrix}2x&1\\y&x\end{pmatrix}
 =\begin{pmatrix}2x+2xy^2&1+2x^2y\end{pmatrix}.
\]
Jest to ten sam wynik, który otrzymujemy przez
różniczkowanie funkcji $x^2+y+x^2y^2$.
\end{przyklad}

\begin{przyklad}[Pochodne częściowe funkcji dwóch zmiennych funkcyjnych]
Niech $X=C([0,1],\R)$ z normą supremum.
Dla $H:X\times X\to X$, $H(u,v)=uv+u^2$,
mamy $D_1H(u,v)h=(v+2u)h$ oraz
$D_2H(u,v)k=uk$. Istotnie, przyrost pierwszego
bloku daje resztę $h^2$ o normie $\|h\|_\infty^2$,
a przyrost drugiego bloku jest dokładnie liniowy.
Pochodne częściowe są ciągłe w normie operatorowej:
dla operatora mnożenia $M_w:h\mapsto wh$ zachodzi
$\|M_w\|\le\|w\|_\infty$, więc różnica operatorów
jest oszacowana przez normę różnicy ich mnożników.
Z twierdzenia \ref{ma2:thm:continuous-block-partials}
otrzymujemy $DH(u,v)(h,k)=(v+2u)h+uk$.
\end{przyklad}

\section{Interpretacja geometryczna}

Dla $f:U\subseteq\R^n\to\R$ różniczkowalnej
w $p$ afiniczny wykres przybliżenia jest zbiorem
\[
 \{(p+h,f(p)+Df(p)h):h\in\R^n\}.
\]
Jego odchylenie pionowe od wykresu $f$ wynosi
$r(h)=o(\|h\|)$ na mocy definicji
\ref{ma2:def:frechet}. Dla $n=2$ jest to płaszczyzna
styczna w punkcie $(p,f(p))$, a dla dowolnego $n$
hiperpłaszczyzna styczna. Nie wnioskujemy jeszcze,
że cały wykres jest gładką podrozmaitością:
do tego potrzebna jest odpowiednia regularność w otoczeniu.

\begin{przyklad}
Dla $f(x,y)=x^2+2y^2$ w $(1,1)$ otrzymujemy
$\mathrm df_{(1,1)}(h,k)=2h+4k$.
Płaszczyzna styczna do wykresu ma równanie
$z=3+2(x-1)+4(y-1)$.
Po podstawieniu $x=1+h$, $y=1+k$ różnica wysokości
wynosi $h^2+2k^2$, co nie przekracza $2(h^2+k^2)$.
Po podzieleniu przez $\sqrt{h^2+k^2}$ iloraz dąży
do zera. Jest to bezpośrednie sprawdzenie
liniowego przybliżenia, a nie tylko obliczenie
pochodnych cząstkowych.
\end{przyklad}

W następnym rozdziale wykorzystamy regułę łańcuchową
i lemat o przyroście do ogólnych nierówności o wartości
średniej. Wzór całkowy dla krzywych o wartościach
w przestrzeni Banacha wymaga osobnego uzasadnienia
istnienia całki.
